% q0032.tex — So where is this discount factor, anyway?
\question{
  id        = {0032},
  title     = {So where is this discount factor, anyway?},
  source    = {OBECON},
  year      = {2026},
  round     = {Seletiva IEO -- Phase B},
  language  = {English},
  topic     = {Finance},
  subtopic  = {Asset Pricing / CAPM / Fama-French},
  type      = {dissertative},
  statement = {%
    Asset pricing theory provides a framework for determining the fair value
    of financial assets based on expected cash flows and associated risks.

    \begin{enumerate}[(a)]
      \item Consider a risk-free asset paying cash flow $C$ at $t+1$, priced
            today as $P_t = C/(1+r)$. Explain the intuition behind this
            present value formula. Why are future cash flows discounted, and
            what does the discount rate represent?

      \item When investing in stocks, what are the two main sources of return?

      \item The Capital Asset Pricing Model (CAPM) gives the required return
            on stock~$i$ as
            \[
              R_i^e = R_f + \beta_i(R_m - R_f).
            \]
            What does $(R_m - R_f)$ represent? Why should stocks with higher
            $\beta$ require higher expected returns?

      \item The Fama--French three-factor model extends CAPM:
            \[
              R_i^e = R_f + \beta_i(R_m - R_f) + s_i \cdot SMB + h_i \cdot HML,
            \]
            where $SMB > 0$ and $HML > 0$ empirically. Formulate one
            hypothesis explaining why small firms outperform large firms, and
            one hypothesis explaining why value firms outperform growth firms.
            \textit{Hint: consider how value and growth firms perform during
            economic downturns.}
    \end{enumerate}
  },
  choices   = {},
  answer    = {},
  solution  = {%
    \textbf{(a)} Money has time value: a dollar today can be invested to
    earn returns, so future cash flows must be discounted to reflect the
    opportunity cost of waiting. The discount rate $r$ represents the
    required rate of return --- the compensation investors demand for
    deferring consumption. For risk-free assets this equals the risk-free
    rate.

    \textbf{(b)} Dividends (periodic cash distributions) and capital gains
    (appreciation in share price realised upon sale).

    \textbf{(c)} $(R_m - R_f)$ is the \emph{market risk premium}: the extra
    return investors demand for bearing systematic market risk. $\beta_i$
    measures a stock's exposure to systematic risk; since investors are
    risk-averse, higher-beta stocks must offer higher expected returns to
    compensate for greater risk.

    \textbf{(d)} \emph{Size premium}: small firms face greater credit,
    liquidity, and operational risks than large firms, requiring higher
    returns as compensation. \emph{Value premium}: value firms have
    short-duration cash flows concentrated in the near term; in recessions
    these near-term cash flows deteriorate sharply, exposing value investors
    to higher systematic (cyclical) risk. Investors demand a premium for
    bearing this downturn risk.
  },
}
